\subsection{ 5 мая. }


\begin{definition}
    Сверка в $L_1(\R)$: $f, g \in L_1(\R)$
    \[ (f * g)(x) = \int_{y=-\infty}^{\infty} f(y)g(x-y)dy\]
\end{definition}

\begin{note}
    \textbf{Корректность:} Почему $f * g \in L_1(\R)$?
\end{note}

\begin{proof}
Рассмотрим
\[
\int_x dx \left| \int_y dy f(y) g(x - y) \right|
\]

далее достаточно оценить сверху этот интеграл $\|f\|_{L_1} \cdot \|g\|_{L_1}$ и применить теорему Фубини.
\end{proof}

\begin{statement}
    Свёртка~--- коммутативная операция (доказательство: просто линейная замена)
    \[(f*g)(x) = (g * f)(x).\]
\end{statement}

\begin{note}
\textbf{Зачем она нам нужна?} В пространстве $L_1, \; f, g \in L_1$ можно рассматривать $f + g$ и $f \cdot g$, но последнее может быть не интегрируемо. Поэтому можно рассматривать $f * g$ как <<сверточное умножение>>~--- оно будет в $L_1$.
\end{note}

\begin{example}
    Рассмотрим частные суммы ряда Фурье \[S_n(f; x) = \frac{1}{\pi} \int_{-\pi}^{\pi}D_n(x-t)f(t)dt.\]
    Из этой формулы потом много всего выводилось.
\end{example}

\begin{theorem}
    $f,g \in L_1$. Тогда
    \[ \widehat{f * g} = \widehat{f} \cdot \widehat{g}.\]
\end{theorem}
\begin{proof}
    \[ \widehat{(f * g)}(z) = \int_x dx \left(\int_y dy \; f(y) g(x-y)\right) \cdot e^{-ixz} =\]
    Запускаем теорему Фубини. При взятии модуля экспонента пропадет, значит все остальное можно переставлять с точки зрения теоремы Фубини
    \[ = \int dy \int dx \; (f(y) g(x-y) e^{-ixz}) = \]
    Сделаем замену $t = x-y$
    \[ = \int dy f(y) \int dt g(t) e^{-i(t+y)z} = \int f(y) \widehat{g}(z) e^{iyz} dy = \widehat{f}(z) \cdot \widehat{g}(z) \]

\end{proof}

\begin{example}[Свертка вещь естественная]
    Дискретный случай: 
    \[ (\sum a_k z^k)(\sum b_m z^m) = \sum_n (\sum_k a(k) b(n-k)) z^n\]
\end{example}

\begin{example}
    Комплексная форма записи ряда Фурье может быть полезной \[\sum c_n (e^{ix})^n.\]
\end{example}

\subsubsection{Зачем нужна свертка}

Имеем  $f \in S', \varphi \in S$
\[
<\widehat{f}, \varphi> = \int \widehat{f}(x) \varphi(x) dx = \int \widehat{\varphi]}(x) f(x) dx
\]
Тогда \(<\widehat{f}, \varphi> =_{\textit{def}} <f, \widehat{\varphi}>\) и $\delta * \varphi = \varphi$.

Если мы выберем последовательность регулярных функци $\delta_n \to^{S'} \delta$, то получим:
\[ \delta_n * \varphi \to \delta * \varphi = \varphi \]
То есть можно получить последовательность хорошего вида, сходящуюся к $\varphi$

\begin{note}[Подарок Дня]
    \begin{align*}
        &\forall f \in C[0, 1] \\
        &P_n(x) = \sum_{k = 0}^n C^k_n c^k (1-x)^{n - k} \cdot f\left(\frac{k}{n}\right)\\
        &P_n(x) \uniconv f
    \end{align*}
\end{note}

\begin{example}
    $f(x)\geq 0, \; L\int f(x)dx=1$

    \[ \frac{1}{\varepsilon} f(\frac{x}{\varepsilon} \to_{D'} \delta, \; \varepsilon \to 0\]
    \[ \int_{-\infty}^{\infty} \frac{1}{\varepsilon} f(\frac{x}{\varepsilon}) \varphi(x) dx = \int_{-\infty}^{\infty} f(y) \varphi(\varepsilon y) dy\]

    \( f(y)\varphi(\varepsilon y) = g_\varepsilon (y) \). Тогда по теореме Лебега (в силу финитности $\varphi$)

    \[ \lim_{\varepsilon \to 0+} \int_{-\infty}^{\infty} f(y) \varphi(\varepsilon y) dy = \varphi(0) = <\delta, \varphi>\]
\end{example}
